% Generated draft: requires mathematical and source review.
\documentclass[11pt]{article}
\usepackage[margin=25mm]{geometry}
\usepackage{fontspec}
\setmainfont{Latin Modern Roman}
\usepackage{amsmath,amssymb,mathtools,unicode-math}
\setmathfont{Latin Modern Math}
\usepackage{enumitem,needspace,xurl,hyperref}
\setlength{\parindent}{0pt}
\setlength{\parskip}{5pt}
\setlength{\emergencystretch}{3em}
\newcommand{\N}{\mathbb{N}}
\newcommand{\Z}{\mathbb{Z}}
\newcommand{\Q}{\mathbb{Q}}
\newcommand{\R}{\mathbb{R}}
\newcommand{\C}{\mathbb{C}}
\newcommand{\F}{\mathbb{F}}
\newcommand{\A}{\mathbb{A}}
\newcommand{\PP}{\mathbb{P}}
\newcommand{\E}{\mathbb{E}}
\newcommand{\eps}{\varepsilon}
\newcommand{\dd}{\,\mathrm{d}}
\begin{document}
\section*{UnsolvedMath ID: 8800048}
\textbf{Source problem code:} catalogue entry
\subsection{Definitions and mathematical statement}
A principally polarized abelian threefold over a field $k$ is a pair $(A,\lambda)$ where $A/k$ is an abelian variety of dimension three and $\lambda:A\to A^\vee$ is an ample-line-bundle polarization that is an isomorphism. It is a Jacobian over $k$ if there exists a smooth projective geometrically connected genus-three curve $C/k$ and an isomorphism of principally polarized varieties
\[
(A,\lambda)\cong(\operatorname{Jac}(C),\lambda_C).
\]
This differs from becoming a Jacobian only over $\overline k$ or merely being isogenous to one.

Give an effective input representation of $(A,\lambda)$ suitable for deciding this Jacobian property. The representation should include the principal polarization, not just the underlying abelian variety, and provide finite accessible data for whatever arithmetic/geometric operations the intended decision procedure requires.

The source does not specify the base field, its computational presentation, characteristic restrictions, or admissible exact/approximate input model. Algebraic equations with group/polarization data, period matrices, and theta/level structures are possible but inequivalent frameworks; the summary mentions them without selecting one. It also does not state whether the target is the base-field or geometric Jacobian property. Consequently the task is to define a representation/decision framework for a specified category rather than to assume a uniform finite encoding of all complex data or an unprovided general Schottky test.
\subsection{Short English statement}
Specify effective data for a principally polarized abelian variety of dimension three that lets an algorithm decide whether it is a genus-three Jacobian. The base field, polarization encoding, and descent versus geometric question must be fixed.
\subsection{Sources}
\begin{itemize}
\item Everett Howe, IPAM \emph{Workshop on Number Theory and Cryptography --- Open Problems} (2006), attachment 6237, PDF pages 33--34, schedule source \url{https://www.ipam.ucla.edu/programs/workshops/workshop-i-number-theory-and-cryptography-open-problems/?tab=schedule}; accessed 2026-07-29. The supplied provenance records a PDF/legacy-PPT visual audit.
\item The supplied summary names Serre's obstruction $\chi_{18}$ and period/theta/normalized representations without a full reference, computational hypotheses, or a universally applicable criterion.
\item Local registry \texttt{data/unsolvedmath/problems.json}, permanent UnsolvedMath ID $8800048$, source code catalogue entry, \emph{Effective representation of principally polarized abelian threefolds}. Dataset attribution: Ulam AI and the UnsolvedMath Contributors.
\end{itemize}
\end{document}
